\scanpage{8}
Alors $R$ est muni d'une structure modérée canonique (en tant que modèle $\subset R^1=R$), i.e.
\[
 R\in\operatorname{Ob}\mathcal E\qquad(\mathcal E=\operatorname{Esp}_{\mathcal M_*})
\]
et $\forall$ objet de $\mathcal E$ est isomorphe à un sous-objet d'un $R^n$ \inserted{(produit dans $\mathcal E$)}. On récupère $\mathcal M_n$ comme l'ens des \inserted{sous-ens. de $R^n$ \unsure{correspondants dans Ens}} sous-objets (dans $\mathcal E$) de l'objet $R^n$. On trouve que la donnée d'un $\mathcal M_*$ est équivalente à la donnée d'un système (à isom. près)
\[
 (\mathcal E,R,X\mathrel{\mapstochar\relbar}|X|)
\]
de a) la catégorie $\mathcal E$ avec $\varprojlim$ finies

b) le foncteur $X\mathrel{\mapstochar\relbar}|X|$ $\mathcal E\mathbin{-}(\mathrm{Ens})$ exact à gauche

c) l'objet $R\in\operatorname{Ob}\mathcal E$ tel que tt objet de $R$ se plonge dans un $R^n$.

\margin{Parler aussi des applications modérées d'ens. modérés \uline{dans des $R^I$}.}

\margin{Introduire les parties modérées d'un espace modéré correspondant aux sous-objets dans $\operatorname{Esp}_{\mathcal M_*}$, les propriétés de stabilité $f(A)$, $f^{-1}(B)$, produits cartésiens.}

\figure{\begin{tikzpicture}\node[draw,ellipse] {III};\end{tikzpicture}}
\uline{Limites inductives finies d'espaces modérés}

S'intéresser à la question des réunions finies des parties modérées, et des $\varinjlim$ finies dans $(\operatorname{Esp}_{\mathcal M_*})$.

\margin{A partir d'ici, on suppose que $\unsure{M4|M5}$ de réunion seulement \unsure{possible}. Je verrai \uline{bien} $(M1)_0$ pour la suite par un axiome.}
\mspage{7}
